\chapter{Funding, Margin and Capital Adjustments}\label{ch:rc:funding-margin-and-capital-adjustments}

In 2012 John Hull and Alan White asked in print whether funding is a cost for derivatives
desks, and argued that it should not enter a derivative's price. A practitioner's reply
appeared the same year under the title ``Yes, FVA is a cost for derivatives desks''. In January
2014 JPMorgan reported a 1.5 billion dollar loss from implementing a \omterm{def:rc:funding-margin-and-capital-adjustments:fva}{funding valuation adjustment},
describing an industry migration towards pricing the cost of funding. Both sides were right about something: the debate turned on whose
value is being measured, the firm's as a whole or its shareholders', and on what an
uncollateralised trade forces a bank to borrow. This chapter builds the adjustments that
followed CVA and DVA: the cost of funding uncollateralised positions, of posting initial
margin, and of holding capital against counterparty and CVA risk, and stacks them into the
price a bank quotes.

\section{Funding a derivative: the debate}

A bank that trades an uncollateralised swap with a client hedges it with an opposite swap
under a CSA with the market. When the client's swap is in the bank's favour, the hedge is
against it: the bank posts variation margin on the hedge and receives nothing from the
client. It must borrow that cash at its own funding rate, above the overnight rate paid on
collateral. When the client's swap is against the bank, the bank receives margin on the
hedge and can lend it or reduce its borrowing
(\cref{fig:rc:funding-margin-and-capital-adjustments:flows}).

\begin{omfigure}
\begin{tikzpicture}[font=\small,box/.style={draw,rounded corners,minimum width=2.4cm,minimum height=0.9cm,align=center}]
\node[box] (c) at (0,0) {client\\(no CSA)};
\node[box,fill=omDef!10] (b) at (5,0) {bank};
\node[box] (h) at (10.4,0) {hedge\\counterparty (CSA)};
\node[box,fill=omThm!10] (t) at (5,-2.8) {bank treasury\\(borrows at $r+s_f$)};
\draw[<->,thick] (c) -- node[above] {swap} (b);
\draw[<->,thick] (b) -- node[above] {opposite swap} (h);
\draw[->,omThm,thick] (b.south east) to[bend right=25] node[below,align=center] {variation margin posted\\when the client owes the bank} (h.south west);
\draw[->,omThm,thick] (t) -- node[left,align=right] {cash\\borrowed} (b);
\end{tikzpicture}

\omcaption{Why an uncollateralised trade needs funding: the bank hedges it with a
collateralised trade, posts margin on the hedge when the client's trade is in its favour,
and borrows that margin from its treasury at its funding spread. Schematic.}\label{fig:rc:funding-margin-and-capital-adjustments:flows}
\end{omfigure}

\begin{definition}[Funding valuation adjustment]\label{def:rc:funding-margin-and-capital-adjustments:fva}
The \emph{funding valuation adjustment}\index{funding valuation adjustment} (FVA) is the
value of the cost of funding the collateral that uncollateralised trades and their hedges
require, less the benefit of the collateral they release:
$\mathrm{FVA} = \mathrm{FCA}-\mathrm{FBA}$.
\end{definition}

\begin{definition}[Funding cost and benefit adjustments]\label{def:rc:funding-margin-and-capital-adjustments:fca}
The \emph{funding cost adjustment}\index{funding cost adjustment} is
$\mathrm{FCA} = \int_0^Ts_f(t)\,\mathrm{DEE}(t)\,Q_C(t)Q_B(t)\,dt$, the funding spread $s_f$ paid on
the discounted \omterm{def:rc:counterparty-exposure:profiles}{expected positive exposure} while both parties survive; the \emph{funding
benefit adjustment}\index{funding benefit adjustment} FBA is the same integral on the
discounted \omterm{def:rc:counterparty-exposure:profiles}{expected negative exposure}, at the rate the bank saves.
\end{definition}

The FBA overlaps with DVA: both are value the bank captures from its own credit spread
on what it owes. Firm-value arguments (Hull and White) say FVA should not affect prices;
shareholder-value arguments (Andersen, Duffie and Song) say funding costs transfer value
to creditors and a desk must charge them. Conventions differ between banks: some count
FCA and FBA symmetrically, some only the cost, some remove the overlap with DVA.

\section{Funding cost and benefit adjustments}

\begin{example}[FVA of the ten-year swap]\label{ex:rc:funding-margin-and-capital-adjustments:fva}
Take chapter~17's ten-year USD~100 million swap on which the bank receives the par rate
of 3.80\% from the BBB counterparty of chapter~18, and a bank funding at 80 basis points
over the overnight rate. With the survival of both parties, the FCA is USD~127\,400 and the
FBA USD~185\,797: the swap is expected to be a liability more than an asset, so the net FVA
is a benefit of USD~58\,397. The FCA alone, a charge many banks book, is 1.5 basis points a
year on the swap. Under a CSA with daily margin the funding exposure almost disappears: the
FCA falls to USD~3\,102.
\end{example}

\omcode{code/firm/xvafund/firm_xvafund.py}{26}{34}{Funding cost and benefit adjustments from discounted exposure profiles and joint survival.}\label{lst:rc:funding-margin-and-capital-adjustments:fva}

\section{Initial margin and MVA}

Uncleared margin rules (One Quant Book~2, chapter~10) and clearing houses require
initial margin, posted by both sides and segregated: it cannot be reused, and the poster
earns little on it while funding it at its own rate.

\begin{definition}[Margin valuation adjustment]\label{def:rc:funding-margin-and-capital-adjustments:mva}
The \emph{margin valuation adjustment}\index{margin valuation adjustment} (MVA) is the value
of the cost of funding the initial margin a position requires over its life:
$\mathrm{MVA} = \int_0^T(s_f-r_{\mathrm{IM}})\,\E[\mathrm{IM}(t)]\,P(0,t)\,dt$.
\end{definition}

\begin{example}[MVA of the swap]\label{ex:rc:funding-margin-and-capital-adjustments:mva}
With Book~2's sensitivity model, initial margin is 2.326 standard deviations of the swap's
value over ten days at a rate volatility of 7 basis points a day: USD~4.22 million today,
falling as the swap's DV01 runs off (\cref{fig:rc:funding-margin-and-capital-adjustments:im}).
Funding it at 80 basis points costs USD~175\,217 over the swap's life.
\end{example}

\begin{omfigure}
\begin{tikzpicture}
\begin{axis}[width=11.5cm,height=5.2cm,
  xlabel={time (years)},ylabel={initial margin (USD million)},
  tick label style={font=\small},label style={font=\small},
  xmin=0,xmax=10,ymin=0,ymax=4.6,grid=major,grid style={gray!25},
  table/col sep=comma]
\addplot[omDef,very thick,const plot] table[x=t,y=im]{figdata/rates-credit-risk/19-funding-margin-and-capital-adjustments/im.csv};
\end{axis}
\end{tikzpicture}

\omcaption{Expected initial margin of the ten-year swap under a sensitivity-based model: it
starts at USD~4.22 million and falls with the swap's remaining DV01. Its funding cost over the
life is the MVA. Data: the chapter's tutorial.}\label{fig:rc:funding-margin-and-capital-adjustments:im}
\end{omfigure}

\section{Capital: counterparty and CVA capital, and KVA}

A bank holds capital against the default of its counterparties and against the volatility
of its CVA. Regulatory formulas turn \omterm{def:rc:counterparty-exposure:netting}{netting sets} into exposures and exposures into capital.

\begin{definition}[Exposure at default, SA-CCR]\label{def:rc:funding-margin-and-capital-adjustments:saccr}
The \emph{exposure at default}\index{exposure at default} (EAD) of a \omterm{def:rc:counterparty-exposure:netting}{netting set} is the
exposure a capital rule charges for its counterparty's default. The \emph{standardised
approach for counterparty credit risk}\index{standardised approach for counterparty credit risk}
(SA-CCR) sets $\mathrm{EAD} = 1.4\,(\mathrm{RC}+\mathrm{multiplier}\times\mathrm{AddOn})$: a replacement cost
($\max(V-C,0)$, or $\max(V-C,\mathrm{Th}+\mathrm{MTA}-\mathrm{NICA},0)$ under margin) plus supervisory add-ons
per asset class (for interest rates, 0.5\% of the notional times a supervisory duration and a
maturity factor), scaled down for excess collateral.
\end{definition}

\omcode{code/firm/xvafund/firm_xvafund.py}{61}{89}{SA-CCR: the interest rate add-on with its maturity buckets, the multiplier and the exposure at default.}\label{lst:rc:funding-margin-and-capital-adjustments:saccr}

\begin{example}[The swap's EAD]\label{ex:rc:funding-margin-and-capital-adjustments:ead}
Today the swap has no replacement cost; its supervisory duration is 7.87 years and its add-on
0.5\% of USD~100 million times 7.87, USD~3.93 million, so its EAD is USD~5.51 million. Under a
CSA with a ten-day margin period of risk the maturity factor is $1.5\sqrt{10/250} = 0.30$ and
the replacement cost floor the MTA of USD~0.5 million: EAD USD~2.35 million. The implementation
reproduces the Basel Committee's worked example (an EAD of 569 thousand on three trades).
\end{example}

\begin{definition}[CVA risk capital]\label{def:rc:funding-margin-and-capital-adjustments:cvacap}
\emph{CVA risk capital}\index{CVA risk capital} is the capital a bank holds against losses
from changes in CVA. In the reduced basic approach it is $0.65\sqrt{(\rho\sum_c\mathrm{SCVA}_c)^2+(1-\rho^2)\sum_c\mathrm{SCVA}_c^2}$
with $\rho = 50\%$ and, per counterparty, $\mathrm{SCVA} = \mathrm{RW}\cdot M\cdot\mathrm{EAD}\cdot\mathrm{DF}/1.4$: a risk weight by sector and
credit quality (3\% for an investment-grade industrial), the effective maturity and a
supervisory discount factor.
\end{definition}

\begin{definition}[Capital valuation adjustment, hurdle rate]\label{def:rc:funding-margin-and-capital-adjustments:kva}
The \emph{capital valuation adjustment}\index{capital valuation adjustment} (KVA) is the value
of the return shareholders require on the capital a trade consumes over its life,
$\mathrm{KVA} = h\int_0^T\E[K(t)]P(0,t)Q_C(t)Q_B(t)\,dt$, where $h$ is the \emph{hurdle
rate}\index{hurdle rate}: the return on capital, above the risk-free rate, that the bank
requires of its activities.
\end{definition}

\begin{example}[KVA of the swap]\label{ex:rc:funding-margin-and-capital-adjustments:kva}
With an illustrative 100\% risk weight and 8\% capital ratio, counterparty capital today is
USD~440\,686; CVA capital, at a 3\% risk weight and an effective maturity of 5.5 years (the
average time of the swap's payments), is USD~368\,944.
Projected along the swap's life with the replacement cost set to the \omterm{def:rc:counterparty-exposure:profiles}{expected exposure}, the
capital peaks at USD~1.18 million after 1.25 years
(\cref{fig:rc:funding-margin-and-capital-adjustments:capital}). At a 10\% \omterm{def:rc:funding-margin-and-capital-adjustments:kva}{hurdle rate} the KVA is
USD~558\,718, the largest adjustment of all.
\end{example}

\begin{omfigure}
\begin{tikzpicture}
\begin{axis}[width=11.5cm,height=5.2cm,
  xlabel={time (years)},ylabel={capital (USD thousand)},
  tick label style={font=\small},label style={font=\small},
  xmin=0,xmax=10,ymin=0,ymax=1700,grid=major,grid style={gray!25},
  legend columns=3,legend style={font=\footnotesize,at={(0.5,-0.25)},anchor=north,draw=none,fill=none,/tikz/every even column/.append style={column sep=8pt}},
  table/col sep=comma]
\addplot[omDef,very thick] table[x=t,y=ccr]{figdata/rates-credit-risk/19-funding-margin-and-capital-adjustments/capital.csv};
\addplot[omMeth,very thick] table[x=t,y=cva]{figdata/rates-credit-risk/19-funding-margin-and-capital-adjustments/capital.csv};
\addplot[omThm,very thick,dashed] table[x=t,y=total]{figdata/rates-credit-risk/19-funding-margin-and-capital-adjustments/capital.csv};
\legend{counterparty credit capital,CVA capital,total}
\end{axis}
\end{tikzpicture}

\omcaption{Projected capital of the uncollateralised ten-year swap: counterparty credit
capital from SA-CCR and CVA capital from the reduced basic approach. Capital rises as the
expected replacement cost builds, then falls with the remaining maturity. Data: the chapter's
tutorial.}\label{fig:rc:funding-margin-and-capital-adjustments:capital}
\end{omfigure}

\begin{dated}{2026-09}{Counterparty and CVA capital}\label{dat:rc:funding-margin-and-capital-adjustments:capital}
The Basel Committee's \omterm{def:rc:funding-margin-and-capital-adjustments:saccr}{standardised approach for counterparty credit risk} dates from March
2014; its revised CVA framework, with the basic and standardised approaches, from July 2020,
as part of the Basel~III standards finalised in December 2017. The Basel Framework's credit
risk standard carries an effective date of 1 January 2023 and a next version from 1 January
2028. The US agencies made the standardised approach mandatory for advanced-approaches banks from
1 January 2022; in the European Union it applies since 28 June 2021 and the revised CVA framework since
1 January 2025. Chapter~23 tracks the market-risk dates.
\end{dated}

\section{The whole stack}

\begin{example}[Pricing the swap]\label{ex:rc:funding-margin-and-capital-adjustments:stack}
Uncollateralised, the swap's adjustments are: CVA USD~272\,251 (on first-to-default), DVA
210\,351, FCA 127\,400, FBA 185\,797, no MVA, KVA 558\,718; in total USD~562\,220, or 6.8 basis
points a year on the swap's fixed rate. Under a CSA with bilateral initial margin: CVA 33\,814,
DVA 17\,745, FCA 3\,102, FBA 3\,367, MVA 175\,217, KVA 156\,573; in total USD~347\,595, 4.2 basis
points (\cref{fig:rc:funding-margin-and-capital-adjustments:stack}). Collateral replaces
credit and funding costs with margin costs, and cuts capital.
\end{example}

\begin{omfigure}
\begin{tikzpicture}
\begin{axis}[width=11.5cm,height=5.4cm,ybar,bar width=10pt,
  symbolic x coords={CVA,-DVA,FCA,-FBA,MVA,KVA,total},xtick=data,
  ylabel={USD thousand},
  tick label style={font=\small},label style={font=\small},
  ymin=-300,ymax=800,grid=major,grid style={gray!25},
  legend columns=2,legend style={font=\footnotesize,at={(0.5,-0.2)},anchor=north,draw=none,fill=none,/tikz/every even column/.append style={column sep=8pt}},
  table/col sep=comma]
\addplot[fill=omThm!60,draw=omThm] table[x=item,y=uncoll]{figdata/rates-credit-risk/19-funding-margin-and-capital-adjustments/stack.csv};
\addplot[fill=omDef!60,draw=omDef] table[x=item,y=csa]{figdata/rates-credit-risk/19-funding-margin-and-capital-adjustments/stack.csv};
\legend{no CSA,CSA and initial margin}
\end{axis}
\end{tikzpicture}

\omcaption{The adjustment stack of the ten-year swap, with signs as they enter the price
charged to the client (DVA and FBA reduce it). Without collateral, KVA and CVA dominate;
with a CSA and initial margin, MVA replaces most of the credit and funding charges. Data:
the chapter's tutorial.}\label{fig:rc:funding-margin-and-capital-adjustments:stack}
\end{omfigure}

The adjustments overlap and interact: DVA and FBA both monetise the bank's own spread;
capital depends on CVA hedges, which reduce CVA capital but need funding; margin reduces CVA
and FVA but creates MVA. A desk computes them on common scenarios and quotes their sum, the
subject of chapter~20.

\section{Tutorial: the adjustment stack}\label{tut:rc:funding-margin-and-capital-adjustments:tutorial}

\begin{tutorial}
\textbf{Goal.} Compute FCA, FBA, MVA and KVA for the ten-year swap, with and without a CSA, and
the whole stack. \textbf{End state:} the numbers of
\cref{ex:rc:funding-margin-and-capital-adjustments:fva,ex:rc:funding-margin-and-capital-adjustments:mva,ex:rc:funding-margin-and-capital-adjustments:ead,ex:rc:funding-margin-and-capital-adjustments:kva,ex:rc:funding-margin-and-capital-adjustments:stack}
and the charts.
\begin{tutsteps}
\item \textbf{Exposure}: chapter~17's swap values; lag~1 for default, lag~0 for funding.
\item \textbf{Funding and margin}: \texttt{fca}, \texttt{fba}, \texttt{mva}.
\item \textbf{Capital}: \texttt{ir\_addon}, \texttt{sa\_ccr\_ead}, \texttt{scva},
\texttt{ba\_cva\_capital}; \texttt{capital\_profile()}; \texttt{kva}.
\item \textbf{Stack}: \texttt{stack\_table()}; \texttt{fig\_rc\_xvafund.py} writes the charts.
\end{tutsteps}
\textbf{What to change next.} Buy CVA hedges and recompute CVA capital with the full basic
approach; raise the \omterm{def:rc:funding-margin-and-capital-adjustments:kva}{hurdle rate} to 12\% and see how the price moves.
\end{tutorial}

\section{Build: funding, margin and capital}\label{bld:rc:funding-margin-and-capital-adjustments:xvafund}

\begin{build}
\bfield{Purpose} The firm's funding, margin and capital adjustments and the regulatory
exposure and CVA capital they use; the inputs of chapter~20's quotes.
\bfield{Interface} \texttt{fca}, \texttt{fba}; \texttt{mva}, \texttt{im\_profile} (over Book~2's
\texttt{firm\_ccpbasis}); \texttt{IrTrade}, \texttt{supervisory\_duration},
\texttt{maturity\_factor}, \texttt{ir\_addon}, \texttt{fx\_addon}, \texttt{multiplier},
\texttt{sa\_ccr\_ead}; \texttt{scva}, \texttt{ba\_cva\_capital}; \texttt{kva}.
\bfield{Rules} SA-CCR for interest rate and FX hedging sets only; reduced basic CVA approach;
risk weights and hurdle as inputs; profiles on chapter~17's grid.
\bfield{Acceptance tests} \texttt{code/firm/xvafund/tests/}: the Basel Committee's example~1
(EAD 569); the margined maturity factor; multiplier bounds; single-counterparty CVA capital;
FCA of a flat profile.
\bfield{Stretch} Credit, equity and commodity add-ons; the standardised approach to CVA
capital from CVA sensitivities; capital projected on each path; a funding curve instead of
a flat spread.
\end{build}

\begin{omsources}
\item C.\ Burgard and M.\ Kjaer, ``Partial differential equation representations of
derivatives with bilateral counterparty risk and funding costs'', \emph{Journal of Credit
Risk} 7(3), 2011, 75--93.
\item A.\ Castagna, ``Yes, FVA is a cost for derivatives desks'', SSRN, 2012 (a reply to Hull
and White).
\item L.\ Andersen, D.\ Duffie and Y.\ Song, ``Funding value adjustments'', \emph{Journal of
Finance} 74(1), 2019, 145--192.
\item A.\ Green, C.\ Kenyon and C.\ Dennis, ``KVA: capital valuation adjustment'', SSRN, 2014.
\item Basel Committee on Banking Supervision, \emph{The standardised approach for measuring
counterparty credit risk exposures}, 2014; \emph{Targeted revisions to the credit valuation
adjustment risk framework}, 2020.
\end{omsources}

\section{Exercises}

\begin{exercise}[$\star$]\label{exo:rc:funding-margin-and-capital-adjustments:1}
A \omterm{def:rc:counterparty-exposure:netting}{netting set} has a flat discounted EE of USD~5 million for four years and the bank funds at
60 basis points. Ignoring survival, what is the FCA?
\end{exercise}

\begin{exercise}[$\star$]\label{exo:rc:funding-margin-and-capital-adjustments:2}
Compute the SA-CCR add-on and EAD of a new unmargined five-year USD~50 million swap.
\end{exercise}

\begin{exercise}[$\star$]\label{exo:rc:funding-margin-and-capital-adjustments:3}
Why does the net FVA of the chapter's swap come out as a benefit?
\end{exercise}

\begin{exercise}[$\star\star$]\label{exo:rc:funding-margin-and-capital-adjustments:4}
Why do FBA and DVA overlap, and how do banks avoid counting both?
\end{exercise}

\begin{exercise}[$\star\star$]\label{exo:rc:funding-margin-and-capital-adjustments:5}
Why is the KVA larger than the CVA for the uncollateralised swap?
\end{exercise}

\begin{exercise}[$\star\star$]\label{exo:rc:funding-margin-and-capital-adjustments:6}
Under the CSA the MVA is the largest item. Why can't the bank avoid it?
\end{exercise}

\begin{exercise}[$\star\star\star$]\label{exo:rc:funding-margin-and-capital-adjustments:7}
\emph{Coding.} Compute the FCA of the swap at funding spreads of 40 and 120 basis points,
with and without the CSA.
\end{exercise}

\begin{exercise}[$\star\star\star$]\label{exo:rc:funding-margin-and-capital-adjustments:8}
\emph{Find the flaw.} ``FVA is our treasury's problem, not the trading desk's; we price at
the overnight rate and let treasury worry about funding.''
\end{exercise}

\section{Problem: The Funding Charge}

\begin{problem}[{Weekend problem --- what a CSA is worth to the bank}]\label{pb:rc:funding-margin-and-capital-adjustments:1}
The bank of the chapter funds at 80 basis points over the overnight rate. A treasurer at the
client asks what the bank would give, in running basis points, for the client to sign a CSA
on its ten-year swap.

\medskip\noindent\textbf{Part I --- Funding.}
\begin{enumerate}
\item Give the FCA and FBA without a CSA, and the net FVA.
\item Give the FCA alone as a running spread.
\item Give the FCA under the CSA and the share of it removed.
\item Why does the funding need depend on the hedge rather than on the client trade?
\item Which of FCA and FBA would a bank that books only funding costs report?
\end{enumerate}

\medskip\noindent\textbf{Part II --- Credit and capital.}
\begin{enumerate}[resume]
\item Give CVA and DVA with and without the CSA.
\item Give the EAD today with and without the CSA.
\item Give the KVA with and without the CSA.
\item Why does the CSA cut KVA by less than CVA?
\item What does initial margin add?
\end{enumerate}

\medskip\noindent\textbf{Part III --- The offer.}
\begin{enumerate}[resume]
\item Give the total adjustment without and with CSA and initial margin.
\item Convert both into running basis points.
\item How much could the bank give the client for signing, in running basis points?
\item What would the client need to post, and what would it cost it?
\item Would a CSA without initial margin be better for both?
\end{enumerate}

\medskip\noindent\textbf{Part IV --- Judgement.}
\begin{enumerate}[resume]
\item Who wins the FVA debate for the pricing of this swap?
\item Why would a bank still charge KVA if capital is not a cash cost?
\item How does the bank's own credit quality enter each adjustment?
\item State the \emph{named result}: the FCA of the uncollateralised swap at 80 basis points
and the share the CSA removes.
\item In one sentence: what does collateral buy a bank?
\end{enumerate}
\end{problem}

\section{Interview questions}

\begin{interviewq}[$\star$ \iqroles{trader, bank}]\label{iq:rc:funding-margin-and-capital-adjustments:1}
Explain FVA with the hedge of an uncollateralised swap.
\end{interviewq}

\begin{interviewq}[$\star\star$ \iqroles{researcher}]\label{iq:rc:funding-margin-and-capital-adjustments:2}
Summarise the FVA debate. Where do you stand?
\end{interviewq}

\begin{interviewq}[$\star\star$ \iqroles{risk, bank}]\label{iq:rc:funding-margin-and-capital-adjustments:3}
Walk through SA-CCR for a single interest rate swap.
\end{interviewq}

\begin{interviewq}[$\star\star$ \iqroles{trader}]\label{iq:rc:funding-margin-and-capital-adjustments:4}
What is MVA and how would you reduce it?
\end{interviewq}

\begin{interviewq}[$\star\star\star$ \iqroles{researcher, bank}]\label{iq:rc:funding-margin-and-capital-adjustments:5}
Why is KVA controversial, and how would you compute it for a new trade?
\end{interviewq}

\begin{interviewq}[$\star\star\star$ \iqroles{developer}]\label{iq:rc:funding-margin-and-capital-adjustments:6}
How would you compute all \omterm{def:rc:credit-and-debit-valuation-adjustments:xva}{XVAs} of a large book consistently and fast?
\end{interviewq}
